% answers_short.tex —— 2016 年·山东卷·文科 · 简短答案速查
% 由 tools/make_short_answers.py 自 answers.tex 的 \gksolhead/\gkanswer 提取，
% 并与 2016/2016_数学真题_山东文_答案要点.md 逐题比对后生成，勿手改（改源 MD 后重跑）。

\begin{gknotice}
\textbf{使用说明}：本节只给最终答案便于快速核对，完整推导见解析版。答案为编写者独立推导，\textbf{非官方标准答案}。
\end{gknotice}

\gkpart{选择题}
\begin{center}\begin{tabular}{ccccccccccc}
\toprule
题号 & 1 & 2 & 3 & 4 & 5 & 6 & 7 & 8 & 9 & 10 \\
\midrule
答案 & A & B & D & C & C & A & B & C & D & A \\
\bottomrule
\end{tabular}\end{center}

\gkpart{填空题}
\begin{itemize}[leftmargin=2.6em]
  \item[\textbf{11.}] \(1\)
  \item[\textbf{12.}] \(\dfrac{4}{3}n(n+1)\)
  \item[\textbf{13.}] \(-5\)
  \item[\textbf{14.}] \(2\)
  \item[\textbf{15.}] \(m\in(3,+\infty)\)
\end{itemize}

\gkpart{解答题}
\begin{itemize}[leftmargin=2.6em]
  \item[\textbf{16.}] (1) \(\dfrac{5}{16}\)（\(xy\leqslant 3\) 的 5 种有序数对：\((1,1)\)、\((1,2)\)、\((1,3)\)、\((2,1)\)、\((3,1)\)）；(2) 获得水杯的概率更大：\(P(\text{水杯})=\dfrac{3}{8}>\dfrac{5}{16}=P(\text{饮料})\)（理由见解析版）
  \item[\textbf{17.}] (1) \(\left[k\pi-\dfrac{\pi}{12},\,k\pi+\dfrac{5\pi}{12}\right]\)（\(k\in\mathbf{Z}\)）；(2) \(g\left(\dfrac{\pi}{6}\right)=\sqrt{3}\)
  \item[\textbf{18.}] (1) 见解析版；(2) 见解析版
  \item[\textbf{19.}] (1) \(b_{n}=3n+1\)；(2) \(T_{n}=3n\cdot 2^{n+2}\)
  \item[\textbf{20.}] (1) \(g(x)=\ln x-2ax+2a\)（\(x>0\)）。当 \(a\leqslant 0\) 时，\(g(x)\) 在 \((0,+\infty)\) 上递增；当 \(a>0\) 时，\(g(x)\) 在 \(\left(0,\dfrac{1}{2a}\right)\) 上递增、在 \(\left(\dfrac{1}{2a},+\infty\right)\) 上递减；(2) \(a\in\left(\dfrac{1}{2},+\infty\right)\)
  \item[\textbf{21.}] (1) \(\dfrac{x^{2}}{4}+\dfrac{y^{2}}{2}=1\)；(2) ① \(\dfrac{k'}{k}=-3\) 为定值（证明见解析版）；② 直线 \(AB\) 的斜率的最小值为 \(\dfrac{\sqrt{6}}{2}\)
\end{itemize}
